% ============================================================
%  CHAPTER 2 — VOCABULARY
% ============================================================
\section{Vocabulary}

Tables are organised by theme. Columns: \textbf{English} | \textbf{Transliteration} | \textbf{Hindi} | \textbf{Notes} (gender, irregular forms, etc.)

% Column types
\newcolumntype{V}{>{\raggedright\arraybackslash}p{2.8cm}}
\newcolumntype{T}{>{\raggedright\arraybackslash}p{2.6cm}}
\newcolumntype{D}{>{\raggedright\arraybackslash}p{2.4cm}}
\newcolumntype{N}{>{\raggedright\arraybackslash}p{3.6cm}}

% ------------------------------------------------------------
\subsection{2.1 Greetings \& Basics}

\begin{tabular}{@{}VTDN@{}}
  \toprule
  English & Transliteration & Hindi & Notes \\
  \midrule
  Hello / Namaste      & namaste         & \hw{नमस्ते}          & universal greeting \\
  How are you?         & aap kaise hain? & \hw{आप कैसे हैं?}   & formal \\
  I'm fine             & main thiik hoon & \hw{मैं ठीक हूँ.}  & \\
  Thank you            & shukriyaa       & \hw{शुक्रिया}        & also \textit{dhanyavaad} \\
  You're welcome       & koi baat nahin & \hw{कोई बात नहीं}   & lit.\ "no matter" \\
  Yes                  & haan           & \hw{हाँ}             & \\
  No                   & nahin          & \hw{नहीं}            & \\
  Please               & kripayaa        & \hw{कृपया}           & formal \\
  Sorry / Excuse me    & maaph kijiye   & \hw{माफ़ कीजिए}     & formal \\
  Goodbye              & alvida          & \hw{अलविदा}          & also \textit{namaste} \\
  My name is \ldots    & meraa naam \ldots hai & \hw{मेरा नाम … है.} & \\
  Nice to meet you     & aapse milkar khushii huii & \hw{आपसे मिलकर ख़ुशी हुई.} & formal \\
  \bottomrule
\end{tabular}

% ------------------------------------------------------------
\subsection{2.2 Numbers (\hw{गिनती})}

\begin{tabular}{@{}VTDN@{}}
  \toprule
  English & Transliteration & Hindi & Notes \\
  \midrule
  0  & shunya  & \hw{शून्य}  & \\
  1  & ek      & \hw{एक}    & \\
  2  & do      & \hw{दो}    & \\
  3  & tiin    & \hw{तीन}   & \\
  4  & chaar   & \hw{चार}   & \\
  5  & paanch & \hw{पाँच}  & \\
  6  & chhe    & \hw{छह}    & \\
  7  & saat    & \hw{सात}   & \\
  8  & aath   & \hw{आठ}    & \\
  9  & nau     & \hw{नौ}    & \\
  10 & das     & \hw{दस}    & \\
  11 & gyaarah    & \hw{ग्यारह}    & \\
  12 & baarah     & \hw{बारह}      & \\
  13 & teerah     & \hw{तेरह}      & \\
  14 & chaudah    & \hw{चौदह}      & \\
  15 & pandrah    & \hw{पंद्रह}    & \\
  16 & solah      & \hw{सोलह}      & \\
  17 & satrah     & \hw{सत्रह}     & \\
  18 & aTHaarah   & \hw{अठारह}     & \\
  19 & unniis     & \hw{उन्नीस}    & \\
  20 & biis       & \hw{बीस}       & \\
  21 & ikkiis     & \hw{इक्कीस}    & \\
  22 & baaiis     & \hw{बाईस}      & \\
  23 & teiis      & \hw{तेईस}      & \\
  24 & chaubiis   & \hw{चौबीस}     & \\
  25 & pacchiis   & \hw{पच्चीस}    & \\
  26 & chhabbiis  & \hw{छब्बीस}    & \\
  27 & sattaaiis  & \hw{सताईस}     & \\
  28 & aTHTHaaiis & \hw{अट्ठाईस}   & \\
  29 & unatiis    & \hw{उनतीस}     & \\
  30 & tiis       & \hw{तीस}       & \\
  50 & pachaas    & \hw{पचास}      & \\
  100 & sau       & \hw{सौ}        & \\
  \bottomrule
\end{tabular}

% ------------------------------------------------------------
\subsection{2.3 Time \& Days}

\begin{tabular}{@{}VTDN@{}}
  \toprule
  English & Transliteration & Hindi & Notes \\
  \midrule
  today      & aaj        & \hw{आज}       & \\
  yesterday / tomorrow & kal & \hw{कल} & same word; context clarifies \\
  day after tomorrow & parso & \hw{परसों} & also means day before yesterday \\
  now        & abhi       & \hw{अभी}      & \\
  later      & baad men   & \hw{बाद में}  & \\
  morning    & subah      & \hw{सुबह}     & f. \\
  afternoon  & dopahar    & \hw{दोपहर}    & f. \\
  evening    & shaam      & \hw{शाम}      & f. \\
  night      & raat       & \hw{रात}      & f. \\
  Monday     & somvaar    & \hw{सोमवार}   & m. \\
  Tuesday    & mangalvaar & \hw{मंगलवार}  & m. \\
  Wednesday  & budhvaar   & \hw{बुधवार}   & m. \\
  Thursday   & guruvaar   & \hw{गुरुवार}  & m. \\
  Friday     & shukravaar & \hw{शुक्रवार} & m. \\
  Saturday   & shanivaar  & \hw{शनिवार}   & m. \\
  Sunday     & ravivaar   & \hw{रविवार}   & m. \\
  \bottomrule
\end{tabular}

% ------------------------------------------------------------
\subsection{2.4 People \& Family}

\begin{tabular}{@{}VTDN@{}}
  \toprule
  English & Transliteration & Hindi & Notes \\
  \midrule
  man / person & aadmii  & \hw{आदमी}   & m. \\
  woman        & aurat   & \hw{औरत}    & f. \\
  boy          & ladkaa & \hw{लड़का}  & m. \\
  girl         & ladkii & \hw{लड़की}  & f. \\
  child        & bachcha & \hw{बच्चा}  & m. (bachchi f.) \\
  mother       & maa     & \hw{माँ}    & f.; formal: \textit{maataa} \\
  father       & pitaa   & \hw{पिता}   & m.; informal: \textit{baap} \\
  brother      & bhaaii  & \hw{भाई}    & m. \\
  sister       & bahan   & \hw{बहन}    & f. \\
  husband      & pati    & \hw{पति}    & m. \\
  wife         & patni   & \hw{पत्नी} & f. \\
  friend       & dost    & \hw{दोस्त}  & m./f. \\
  teacher            & adhyaapak & \hw{अध्यापक} & m. (adhyaapikaa f.) \\
  son                & beTaa     & \hw{बेटा}    & m. \\
  daughter           & beTii     & \hw{बेटी}    & f. \\
  grandfather (pat.) & daadaa    & \hw{दादा}    & m. \\
  grandmother (pat.) & daadii    & \hw{दादी}    & f.; NB: \textit{daaRhii} (\hw{दाढ़ी}) = beard! \\
  \bottomrule
\end{tabular}

% ------------------------------------------------------------
\subsection{2.5 Common Verbs}

\begin{tabular}{@{}VTDN@{}}
  \toprule
  English (infinitive) & Transliteration & Hindi & Notes \\
  \midrule
  to be         & honaa    & \hw{होना}    & irregular; highly important \\
  to have       & —        & —            & use \textit{ke paas honaa} \\
  to go         & jaanaa   & \hw{जाना}   & \\
  to come       & aanaa    & \hw{आना}    & \\
  to eat        & khaanaa  & \hw{खाना}   & \\
  to drink      & piinaa   & \hw{पीना}   & \\
  to do / make  & karnaa   & \hw{करना}   & \\
  to say        & kahnaa   & \hw{कहना}   & \\
  to see        & dekhnaa  & \hw{देखना}  & \\
  to want       & chaahna  & \hw{चाहना}  & \\
  to know       & jaannaa  & \hw{जानना}  & \\
  to give       & denaa    & \hw{देना}   & \\
  to take       & lenaa    & \hw{लेना}   & \\
  to live/stay  & rahnaa   & \hw{रहना}   & \\
  to speak      & bolnaa   & \hw{बोलना}  & \\
  to learn      & sikhnaa  & \hw{सीखना}  & \\
  to read       & padhnaa & \hw{पढ़ना}  & also "to study" \\
  to write      & likhnaa  & \hw{लिखना}  & \\
  to understand & samajhnaa & \hw{समझना} & \\
  to think      & sochnaa   & \hw{सोचना}      & \\
  to teach      & sikhaanaa & \hw{सिखाना}     & cf.\ \textit{sikhnaa} = to learn \\
  to meet       & milnaa    & \hw{मिलना}      & \\
  to cook       & pakaanaa  & \hw{पकाना}      & \\
  to drive/ride & chalaanaa & \hw{चलाना}      & \\
  to wake up    & jaagnaa   & \hw{जागना}      & \\
  to like       & pasand honaa & \hw{पसंद होना} & use with \textit{ko}: \textit{mujhe X pasand hai} \\
  \bottomrule
\end{tabular}

% ------------------------------------------------------------
\subsection{2.6 Adjectives}

\begin{tabular}{@{}VTDN@{}}
  \toprule
  English & Transliteration (m.sg.) & Hindi & Notes \\
  \midrule
  good/nice   & acchaa   & \hw{अच्छा}   & inflects \\
  bad         & buraa    & \hw{बुरा}    & inflects \\
  big         & badaa   & \hw{बड़ा}    & inflects \\
  small       & chhotaa  & \hw{छोटा}    & inflects \\
  new         & nayaa    & \hw{नया}     & inflects \\
  old         & puraanaa & \hw{पुराना}  & inflects (for objects; for people: \textit{buu.Rhaa}) \\
  hot         & garam    & \hw{गरम}     & invariable \\
  cold        & thandaa & \hw{ठंडा}   & inflects \\
  beautiful   & sundar   & \hw{सुंदर}   & invariable \\
  expensive   & mahangaa & \hw{महँगा}  & inflects \\
  cheap       & sastaa  & \hw{सस्ता}   & inflects \\
  correct     & sahii     & \hw{सही}    & invariable \\
  \bottomrule
\end{tabular}

% ------------------------------------------------------------
\subsection{2.7 Food \& Drink}

\begin{tabular}{@{}VTDN@{}}
  \toprule
  English & Transliteration & Hindi & Notes \\
  \midrule
  water  & paanii  & \hw{पानी}  & m. \\
  tea    & chaay   & \hw{चाय}   & f. \\
  food   & khaanaa & \hw{खाना}  & m. \\
  bread  & rotii  & \hw{रोटी}  & f. \\
  rice       & chaaval        & \hw{चावल}          & m. \\
  milk       & duudh          & \hw{दूध}            & m. \\
  yogurt/curd & dahii         & \hw{दही}            & m. \\
  egg        & anDaa          & \hw{अंडा}           & m. \\
  potato     & aaluu          & \hw{आलू}            & m. \\
  apple      & seb            & \hw{सेब}            & m. \\
  pineapple  & ananaas        & \hw{अनानास}         & m. \\
  salt       & namak          & \hw{नमक}            & m. \\
  sugar      & chiinii/shakkar & \hw{चीनी / शक्कर}  & f. \\
  spices     & masaalaa       & \hw{मसाला}          & m. \\
  chili      & mirch          & \hw{मिर्च}          & f. \\
  eggplant   & baigan         & \hw{बैंगन}          & m. \\
  okra       & bhindi         & \hw{भिंडी}          & f. \\
  \bottomrule
\end{tabular}

% ------------------------------------------------------------
\subsection{2.8 Places \& Locations}

\begin{tabular}{@{}VTDN@{}}
  \toprule
  English & Transliteration & Hindi & Notes \\
  \midrule
  house / home  & ghar       & \hw{घर}        & m. \\
  room          & kamraa     & \hw{कमरा}      & m. \\
  shop / store  & dukaan     & \hw{दुकान}     & f. \\
  city / town   & shahar     & \hw{शहर}       & m. \\
  village       & gaanv      & \hw{गाँव}      & m. \\
  mountain      & pahaaR     & \hw{पहाड़}     & m. \\
  garden        & bagiicha   & \hw{बगीचा}     & m. \\
  place         & jagah      & \hw{जगह}       & f.; formal: \textit{sthal} \\
  office        & aafis      & \hw{ऑफिस}      & m. \\
  school        & skuul      & \hw{स्कूल}     & m. \\
  temple        & mandir     & \hw{मंदिर}     & m. \\
  church        & girjaaghar & \hw{गिरजाघर}   & m. \\
  country       & desh       & \hw{देश}       & m. \\
  world         & vishv      & \hw{विश्व}     & m. \\
  \bottomrule
\end{tabular}

% ------------------------------------------------------------
\subsection{2.9 Connectors, Adverbs \& Useful Words}

\begin{tabular}{@{}VTDN@{}}
  \toprule
  English & Transliteration & Hindi & Notes \\
  \midrule
  and              & aur           & \hw{और}             & \\
  or               & yaa           & \hw{या}              & \\
  but              & lekin         & \hw{लेकिन}           & \\
  so / then        & to            & \hw{तो}              & \\
  this is why      & isliye        & \hw{इसलिए}           & \\
  fortunately      & saubhaagya se & \hw{सौभाग्य से}      & \\
  a little         & thoRaa        & \hw{थोड़ा}           & inflects \\
  a lot / more     & bahut/zyaadaa & \hw{बहुत / ज़्यादा}  & \\
  different/various & alag         & \hw{अलग}             & \textit{alag-alag} = various \\
  first            & pehlaa        & \hw{पहला}            & inflects \\
  again            & phir          & \hw{फिर}             & \textit{phir ek baar} = once more \\
  really? / truly? & sach men?     & \hw{सच में?}         & \\
  week             & saptaah       & \hw{सप्ताह}          & m.; \textit{saptaahant} = weekend \\
  hour             & ghantaa       & \hw{घंटा}            & m.; pl.\ \textit{ghante} \\
  question         & savaal        & \hw{सवाल}            & m. \\
  answer           & uttar         & \hw{उत्तर}           & m. \\
  \bottomrule
\end{tabular}
